% !TeX spellcheck = de_DE
\chapter*{Kurzzusammenfassung}



Abgesehen davon, dass Subfaktoren ein interessantes mathematisches Gebiet für sich darstellen, haben sie auch die Aufmerksamkeit von Physikern erregt, da vermutet wird, dass es einen Zusammenhang zwischen Subfaktoren und konformen Feldtheorien (CFT) gibt. Obwohl inzwischen eine überzeugende Menge an Hinweisen für die Gültigkeit dieser Vermutung erbracht wurde, existieren immer noch einige Lücken: Es gibt eine Reihe außergewöhnlicher Subfaktoren, für die keine entsprechende CFT bekannt ist. Daher ist es notwendig, neue Techniken für die Konstruktion einer CFT aus einem Subfaktor zu entwickeln. Hier ist es sinnvoll, die zugrunde liegende mathematische Struktur genauer zu untersuchen: Aus den sogenannten ,,even parts'' eines Subfaktors ergeben sich zwei unitäre Fusionskategorien (UFCs). Ein vielversprechender Ansatz ist, aus diesen Kategorien Quantenspinsysteme zu konstruieren und zu untersuchen, um eine Verbindung zu CFTs zu finden.

Das einfachste Beispiel, das neue Techniken zur Konstruktion einer CFT erfordert, ist der Haagerup Subfaktor, da er der kleinste Subfaktor mit einem Index größer als vier ist. In dieser Arbeit untersuchen wir mithilfe von ein- und zweidimensionalen Gittermodellen die Frage, ob eine CFT existiert, die zum Haagerup Subfaktor gehört. Die erste Aufgabe hierbei besteht darin, die $F$-Symbole der Kategorie zu berechnen, da diese bei allen Modellen, die in dieser Arbeit untersucht werden, einen entscheidenden Bestandteil der Konstruktion darstellen. Wir betrachten die folgenden Modelle:
	\begin{enumerate}
		\item Das sogenannte ,,Golden Chain''-Modell, bei dem es sich um eine eindimensionale Spinkette handelt, deren Grundzustand Informationen (wie zum Beispiel den Wert der zentralen Ladung) über die hypothetische CFT enthält.
		\item Quantenspinketten (wie die Golden Chain) mit Defekten. Die Konstruktion der Vertices, die die Kette mit Defekten bilden, liefert Einblicke in eine mögliche Konstruktion einer unitären modularen Tensorkategorie (UMTC) über die sogenannte Annulare Kategorie.
		\item Das Levin-Wen Modell, welches ein zweidimensionales Gittermodell mit exakt lösbarem Hamiltonian ist, das eine topologische Quantenfeldtheorie liefert. Am interessantesten für uns ist, dass die Anregungen des Systems eine UMTC ergeben.
	\end{enumerate}

Wir stellen fest, dass die Untersuchung der UFCs selbst keine Hinweise auf eine Haagerup CFT liefert, und schließen daraus, dass es notwendig ist, diese Untersuchung auf die entsprechende UMTC auszudehnen, die zum Beispiel über die Anregungen des Levin-Wen Modells konstruiert werden kann.


\hspace{10pt}

\noindent
\textit{Schlagwörter:} Konforme Feldtheorie, Anyon-Ketten, Fusionskategorien